Wir wenden
den [[Mannigfaltigkeit mit Rand/Satz von Stokes/Fakt|Satz von Stokes]]
auf die Kugeln
\mathl{B^n(P, \epsilon)}{} mit dem Rand
\mathl{S^{n-1}(P, \epsilon)}{}
\zusatzklammer {mit Mittelpunkt $P$ und Radius $\epsilon$} {} {} an und erhalten

\mavergleichskettealign
{\vergleichskettealign
{ \operatorname{lim}_{   \epsilon \rightarrow  0  } \, { \frac{   1 }{  \epsilon^n } }     \int_{S^{n-1}(P, \epsilon)} \omega
}
{ =} { \operatorname{lim}_{   \epsilon \rightarrow  0  } \, { \frac{   1 }{  \epsilon^n } }     \int_{ \partial B^{n}(P, \epsilon)}   \omega
}
{ =} { \operatorname{lim}_{   \epsilon \rightarrow  0  } \, { \frac{   1 }{  \epsilon^n } }     \int_{B^{n}(P, \epsilon)} d \omega
}
{ =} { \operatorname{lim}_{   \epsilon \rightarrow  0  } \, { \frac{   1 }{  \epsilon^n } }     \int_{B^{n}(P, \epsilon)} f d \lambda^n
}
{ =} {\beta_n \cdot f(P)
}
}
{}
{}{,}
wobei wir für das letzte Gleichheitszeichen die Stetigkeit von $f$ benutzt haben.

 [[Kategorie:Latexseite]]
